\documentclass[11pt,letterpaper]{article}
\usepackage[utf8]{inputenc}
\usepackage[margin=1in]{geometry}
\usepackage{booktabs}
\usepackage{longtable}
\usepackage{hyperref}
\usepackage{listings}
\usepackage{xcolor}
\usepackage{enumitem}

\definecolor{codegreen}{rgb}{0,0.6,0}
\definecolor{backcolour}{rgb}{0.95,0.95,0.92}
\definecolor{cautionred}{rgb}{0.55,0,0}
\definecolor{okgreen}{rgb}{0,0.45,0}
\definecolor{unverifiedorange}{rgb}{0.75,0.35,0}

\lstset{
    backgroundcolor=\color{backcolour},
    commentstyle=\color{codegreen},
    basicstyle=\ttfamily\footnotesize,
    breaklines=true, captionpos=b, keepspaces=true,
    numbers=none, showspaces=false, showstringspaces=false, tabsize=2
}

\newenvironment{caution}[1]{%
    \par\addvspace{8pt}%
    \noindent{\color{cautionred}\rule{\linewidth}{0.8pt}}\par\vspace{2pt}
    \noindent\textbf{\color{cautionred}Caution: #1}\par\vspace{2pt}
}{%
    \par\vspace{2pt}\noindent{\color{cautionred}\rule{\linewidth}{0.4pt}}%
    \par\addvspace{8pt}%
}

\newcommand{\verified}{{\color{okgreen}\textbf{[VERIFIED]}}}
\newcommand{\unverified}{{\color{unverifiedorange}\textbf{[UNVERIFIED]}}}
\newcommand{\passif}[1]{\par\vspace{2pt}\noindent\textbf{\color{okgreen}Pass:} #1\par\vspace{4pt}}
\newcommand{\failif}[1]{\noindent\textbf{If it fails:} #1\par\vspace{6pt}}

\title{SOP: Bringing Up a NEMA 23 Joint \\ \large Arctos Robotic Arm --- MKS\_57D over CAN}
\author{Applies to Joint X}
\date{September 2026}

\begin{document}
\maketitle

\begin{caution}{This procedure has never been completed on real NEMA 23 hardware}
No NEMA 23 joint on this arm has ever communicated over CAN. The MKS\_57D driver
installed on Joint~X turned out to be the \textbf{RS485 variant}, which has no CAN
transceiver, so every step below that involves the CAN protocol is extrapolated from
the NEMA 17 / Servo42D work rather than confirmed on a 57D.

Each step is tagged \verified{} or \unverified{}:
\begin{itemize}
    \item \verified{} --- confirmed on this arm's hardware, and expected to carry
          over directly (host-side or physical steps).
    \item \unverified{} --- assumed by analogy with the Servo42D. Plausible, but
          nobody has seen it work on a 57D. Confirm each of these as you go and
          update this document with what you actually observe.
\end{itemize}
\end{caution}

\noindent Procedure only --- for protocol details, encoder maths and troubleshooting,
see \texttt{arctos\_motor\_guide.tex}.

\begin{table}[h]
\centering
\begin{tabular}{ll}
\toprule
\textbf{Joint} & X (base/waist rotation) \\
\textbf{CAN ID} & \texttt{0x01} (per config; confirm by scan) \\
\textbf{Gear ratio} & 13.5:1 \\
\textbf{Motor} & NEMA 23, 1.8~N\textperiodcentered m, 2.8~A, 6.35~mm shaft \\
\textbf{Driver} & MKS\_57D \\
\bottomrule
\end{tabular}
\end{table}

\begin{caution}{A NEMA 23 is substantially more dangerous than a NEMA 17}
1.8~N\textperiodcentered m through 13.5:1 is roughly \textbf{24~N\textperiodcentered m
at the joint} --- about 75$\times$ the NEMA 17 joints' output torque, and more than
enough to injure a hand or destroy a printed part before you can react. Keep clear of
the mechanism, keep the emergency stop in a second terminal, and keep a hand on the
power switch for every motion step.
\end{caution}

\section*{Step 0 --- Confirm the board is the CAN variant \verified}
\textbf{Do this first. It is the failure that has already cost this project weeks.}

\begin{enumerate}[leftmargin=*]
    \item Read the part number and silkscreen on the MKS\_57D board.
    \item Confirm it is the \textbf{CAN} variant, not RS485. The two look similar and
          both present a differential pair.
\end{enumerate}

\begin{caution}{RS485 boards will never work, at any bitrate or address}
RS485 carries MKS's own async serial protocol, not CAN framing. A CAN controller and
an RS485 transceiver cannot talk over the same wires. If this board is RS485 you must
either swap it for the CAN variant --- which keeps this entire procedure and every
script valid --- or treat the joint as a separate RS485 subsystem with its own
adapter and its own driver software, which is out of scope here.
\end{caution}
\passif{Board is confirmed CAN variant.}
\failif{Stop. No amount of software work will make an RS485 board appear on the bus.}

\section*{Step 1 --- Pre-power electrical checks \verified}
\begin{enumerate}[leftmargin=*]
    \item Everything powered off, measure \texttt{CAN\_H} to \texttt{CAN\_L}:
          expect $\approx 60\,\Omega$.
    \item Confirm only \texttt{CAN\_H}, \texttt{CAN\_L}, \texttt{GND} run between
          the CANable and the controller.
    \item \textbf{Check the supply.} A NEMA 23 at 2.8~A draws far more than the
          NEMA 17 joints. Confirm the supply is rated for it, and that wiring to the
          driver is sized appropriately.
    \item Confirm the joint is mechanically clear and you know its travel limits.
\end{enumerate}
\passif{$\approx 60\,\Omega$, correct wiring, adequate supply.}

\section*{Step 2 --- Configure the controller panel \unverified}
The 57D's menu layout has not been confirmed to match the 42D's. Set and \textbf{save}:

\begin{table}[h]
\centering
\begin{tabular}{ll}
\toprule
CAN ID & \texttt{0x01} \\
CAN Rate & \texttt{500} \\
CAN Response (CANRSP) & \texttt{ENABLE} \\
CAN CRC & \texttt{SUM-8} \\
Working current & $\leq 2.8$~A (the NEMA 23's rating) \\
\bottomrule
\end{tabular}
\end{table}

Power-cycle after saving.

\begin{caution}{Set working current deliberately}
A 42D in this project was found storing 3000~mA against a 1.2~A motor. On a 2.8~A
NEMA 23 the consequences of an over-set current are correspondingly larger.
\end{caution}
\passif{Settings persist across a power cycle.}
\failif{If the menu differs from the above, record what it actually offers and update
this document --- that is exactly the unverified part.}

\section*{Step 3 --- Bring up the CAN interface \verified}
Identical to the NEMA 17 procedure; this is host-side and motor-independent.
\begin{lstlisting}[language=bash]
lsusb                                    # expect ID 16d0:117e CANable2
ls -l /dev/serial/by-id/
sudo pkill slcand
sudo ip link delete can0 2>/dev/null
sudo slcand -o -c -s6 /dev/serial/by-id/<your-canable> can0
sudo ip link set can0 txqueuelen 1000
sudo ip link set up can0
ip -br link show can0
\end{lstlisting}
\passif{\texttt{can0 \quad UP}}

\section*{Step 4 --- Confirm the real CAN ID \unverified}
The scan itself is verified; what is unverified is whether a 57D answers the
\texttt{0x31} query the same way a 42D does.
\begin{lstlisting}[language=bash]
python3 src/arctos_can_control/arctos_can_control/can_id_scan.py \
    --channel can0 --ids 1-16
\end{lstlisting}
Read-only --- never enables coils, never commands motion.
\passif{Exactly one responder, with a valid checksum.}
\failif{\textbf{Total silence with the board powered and wiring confirmed is the
RS485 signature.} Go back to Step 0 and check the part number before doing anything
else. Note that the slcan adapter does not populate kernel bus-error counters, so
you cannot use ``zero errors'' as corroboration here.}

\section*{Step 5 --- Read-only health check \unverified}
\begin{lstlisting}[language=bash]
python3 src/arctos_can_control/arctos_can_control/joint_reliability_test.py \
    --channel can0 --can-id 0x01 --gear-ratio 13.5 --samples 200
\end{lstlisting}

\begin{table}[h]
\centering
\begin{tabular}{lll}
\toprule
\textbf{Query} & \textbf{Opcode} & \textbf{Expected (by analogy)} \\ \midrule
Encoder          & \texttt{0x31} & stable; 0 jitter, 0 drift \\
Velocity         & \texttt{0x32} & zero \\
Error            & \texttt{0x39} & small \\
Enable state     & \texttt{0x3A} & \texttt{01} \\
Shaft protection & \texttt{0x3E} & \texttt{00} \\
\bottomrule
\end{tabular}
\end{table}

\begin{caution}{Confirm the encoder resolution on the 57D}
This procedure assumes 16{,}384 counts/revolution and 3200 microstep pulses/rev, as
on the 42D. \textbf{Verify this before trusting any angle maths}: command a known
number of pulses and check the encoder delta matches. If the 57D differs, every
conversion in this SOP changes.
\end{caution}
\passif{100\% response, 0 checksum failures, 0 jitter at rest.}

\section*{Step 6 --- Check the mechanism by hand \verified}
\begin{enumerate}[leftmargin=*]
    \item Disable coils: \texttt{cansend can0 001\#F300F4}.
          \emph{(Checksum: \texttt{0xF3+0x00+0x01 = 0xF4}.)}
    \item Confirm \texttt{0x3A} now reads \texttt{00}.
    \item Turn the joint output by hand, gently, both directions, watching the encoder.
\end{enumerate}

\begin{caution}{Do not drop coils on a loaded base joint}
Joint~X carries the entire arm above it. Removing holding torque may let the whole
assembly swing. Only do this with the arm supported or removed.
\end{caution}
\passif{Encoder tracks your hand both ways; motion feels smooth.}
\failif{Gritty resistance means the gear mesh needs plastic-safe grease. No encoder
change while the output turns means the motor is not coupled to it.}

\section*{Step 7 --- First commanded motion \unverified}
Keep the emergency stop in a second terminal before sending anything:
\begin{lstlisting}[language=bash]
cansend can0 001#F7F8
\end{lstlisting}
\emph{(Checksum: \texttt{0xF7+0x01 = 0xF8}.)}

\begin{lstlisting}[language=bash]
python3 src/arctos_can_control/arctos_can_control/joint_unstick_move_test.py \
    --can-id 0x01 --gear-ratio 13.5 \
    --step-degrees 0.25 --num-steps 1 --max-unstick 0 \
    --band-lo -1 --band-hi 1 --speed 25
\end{lstlisting}

\begin{caution}{Start smaller than you would on a NEMA 17}
At 13.5:1 the reduction is far lower than the NEMA 17 joints' 48:1 and 67.82:1, so
the same commanded joint angle turns the motor far less --- and the joint moves much
faster for a given motor speed. A $0.25^\circ$ joint step here is only
$3.4^\circ$ of motor rotation. Verify direction and magnitude on this first step
before increasing anything.
\end{caution}
\passif{$\approx$100\% of commanded, with \texttt{FD=[1, 2]}.}
\failif{Do not retry into a stall. Return to Step 6.}

\section*{Step 8 --- Stepped travel test \unverified}
\begin{lstlisting}[language=bash]
python3 src/arctos_can_control/arctos_can_control/joint_unstick_move_test.py \
    --can-id 0x01 --gear-ratio 13.5 \
    --step-degrees 2.0 --num-steps 5 --max-unstick 0 \
    --band-lo -15 --band-hi 15 --speed 25

# reverse, same step count
python3 src/arctos_can_control/arctos_can_control/joint_unstick_move_test.py \
    --can-id 0x01 --gear-ratio 13.5 \
    --step-degrees -2.0 --num-steps 5 --max-unstick 0 \
    --band-lo -15 --band-hi 15 --speed 25
\end{lstlisting}
\passif{Every step 99--101\% with \texttt{FD 02}; round trip returns to start.}

\section*{Step 9 --- Establish this motor's speed limit \unverified}
\textbf{The NEMA 17 figures do not transfer.} A NEMA 23 has different winding
inductance and a different torque-speed curve, so its usable speed ceiling must be
measured, not assumed.

With the joint unloaded and clear, run short trials at increasing \texttt{speed} and
record actual rpm against commanded. The usable limit is the highest speed where
actual still tracks commanded to within about 1\%; beyond that the motor silently
loses steps.

Record the result here once measured:

\begin{table}[h]
\centering
\begin{tabular}{ll}
\toprule
Linear range (tracks commanded) & \underline{\hspace{3cm}} rpm \\
Roll-off begins                 & \underline{\hspace{3cm}} rpm \\
Saturation ceiling              & \underline{\hspace{3cm}} rpm \\
\bottomrule
\end{tabular}
\caption{To be filled in from measurement. For reference, the NEMA 17 joints track
linearly to 225~rpm and saturate near 310~rpm.}
\end{table}

\section*{Step 10 --- Acceptance}
\begin{table}[h]
\centering
\begin{tabular}{ll}
\toprule
\textbf{Check} & \textbf{Criterion} \\ \midrule
Board variant     & confirmed CAN, not RS485 \\
CAN ID            & confirmed by scan \\
Comms             & 100\% response, 0 checksum failures \\
Encoder at rest   & 0 jitter, 0 drift \\
Encoder scale     & confirmed for the 57D, not assumed \\
Hand check        & smooth, tracks both directions \\
Per-step accuracy & 99--101\% of commanded \\
Status frames     & \texttt{FD 02} on every step \\
Round trip        & returns to start \\
Speed limit       & measured and recorded \\
\bottomrule
\end{tabular}
\end{table}

\section*{Step 11 --- Shutdown \verified}
\begin{lstlisting}[language=bash]
# stop any running script first (Ctrl+C)
cansend can0 001#F300F4     # disable coils, only if safe to do so
sudo ip link set can0 down
sudo pkill slcand
\end{lstlisting}

\section*{Update this document}
Every \unverified{} tag above is a gap. As each step is confirmed on real NEMA 23
hardware, change its tag to \verified{} and record what was actually observed --- in
particular the panel menu layout, the encoder resolution, and the measured speed
limits.

\end{document}
